\section{Actual-field interfaces and curvature interactions}
\label{sec:field-interfaces}

The common problem is the realizability of curvature under actual field
constraints. A numerical curvature budget, a differential compatibility
estimate, and a construction attaining a limiting budget answer different
parts of that problem. The analytical interface can be separated from
monopole supply, while the variational question requires comparing all
curvature costs on the same constructed metric.

\subsection{An analytical interface before the choice of supplier}
For a smooth self-dual two-form $\psi$ set
\[
 q=|\psi|^2,\qquad
 A_\psi=\frac{q}{24}\Id-\frac18\psi\otimes\psi,
 \qquad B_\psi=\fint_M|\Delta_H\psi|^2,
\]
and
\[
 d_\psi=\fint_M\bigl(
 |\Ric+g|^2+|W^+-A_\psi|^2+(q-2)^2\bigr).
\]
Section~\ref{sec:bounded-forms} proves that, on a closed oriented
four-manifold, $\|\psi\|_\infty\le K$ implies
\begin{equation}\label{eq:interfaces-bounded-form}
 (1+B_\psi)d_\psi\ge c(K)>0.
\end{equation}
This statement does not assume a spin structure or a monopole class.
It retains a specific polynomial curvature model, an amplitude bound,
and the actual Hodge-source cost. These are the quantities that a
different field equation would have to supply. In particular, the
fixed model in \eqref{eq:interfaces-bounded-form} is not the zero-Weyl
limit of the parameter theorem.

Ordinary and special monopole classes, stable cohomotopy
\cite{BauerFuruta2004I,Bauer2004II}, and the twisted theory
\cite{IMN2016} provide distinct solvability statements.
The single-solution results of Sections~\ref{sec:tracefree-gap} and
\ref{sec:single-solution-coercivity} use one actual solution on each
relevant negative Yamabe representative. The finite-density theorem
retains its own conformal supply assumptions. Neither the abstract
form criterion nor the name of a supplier identifies these quantifiers.
The finer cancellation in the linear theorem also uses the spinor
curvature action and the determinant-square identity, beyond
\eqref{eq:interfaces-bounded-form} alone.

\subsection{Spectral defects and simultaneous differential constraints}
For a smooth angle $\theta:M\to\mathbb R/(2\pi\mathbb Z)$,
Furuta--Matsuo define \cite[(8)]{FurutaMatsuo2016}
\begin{equation}\label{eq:interfaces-spectral-defect}
 \lambda_\theta(g)=
 \inf_{0\ne\sigma\in H^1(\Lambda^+)}
 \frac{\|(d+d^*)(\sin\theta\,\sigma)\|_2^2
       +\|\nabla(\cos\theta\,\sigma)\|_2^2
       +2\|d|\sigma|\|_2^2}{\|\sigma\|_2^2}.
\end{equation}
Their differential inequality applies to smooth connection--spinor
pairs; their topological curvature conclusions additionally use
perturbed Seiberg--Witten solvability. These are separate interfaces.
In particular, the vanishing of
\eqref{eq:interfaces-spectral-defect} imposes simultaneous differential
conditions, rather than only a scalar energy equality. Comparing those
conditions with the same-field Hessian and Bianchi identities is a
concrete possible extension; no uniform spectral gap across degenerating
metrics is assumed here.

The relation between exact structure and weak-error estimates is also
visible in the compact integrability theorem under the third Gray
curvature condition \cite{ApostolovArmstrongDraghici2002}. Its full
almost-K\"ahler and curvature hypotheses must be retained. Exact
integrability by itself is not a uniform weak-error theorem.
The broader curvature and Einstein background is surveyed in
\cite{LeBrunOverview2026}; the common starting point remains the
classical actual-field and convex-body estimates
\cite{LeBrun2001,LeBrunConvex}.

\subsection{The joint cost of a Weyl-lowering construction}
Write
\[
 \mathcal W_{\mathrm{op}}(g)
   =\int_M\bigl(|W^+|^2+|W^-|^2\bigr),\qquad
 \mathcal W_{\mathrm{ten}}(g)
   =\int_M W_{ijkl}W^{ijkl}=4\mathcal W_{\mathrm{op}}(g).
\]
Here $|W^\pm|$ is the Hilbert--Schmidt norm on two-forms. The
Chern--Gauss--Bonnet identity gives the exact relation
\begin{equation}\label{eq:interfaces-joint-cost}
 P=2\mathcal W_{\mathrm{op}}+\frac{S}{12}-16\pi^2\chi(M)
  =\frac12\mathcal W_{\mathrm{ten}}+\frac{S}{12}-16\pi^2\chi(M).
\end{equation}
Consequently, for two metrics on the same closed smooth manifold,
\begin{equation}\label{eq:interfaces-fixed-topology}
 \Delta P=2\Delta\mathcal W_{\mathrm{op}}+\frac{\Delta S}{12}
         =\frac12\Delta\mathcal W_{\mathrm{ten}}+\frac{\Delta S}{12}.
\end{equation}
Thus $\Delta\mathcal W_{\mathrm{ten}}<0$ lowers $P$ exactly when
$\Delta S<-6\Delta\mathcal W_{\mathrm{ten}}$. A comparison with two
separate summands has a different topological term: if
$Y=M\# Z$ and $\Delta_\# H=H(g_Y)-H(g_M)-H(g_Z)$, then
\begin{equation}\label{eq:interfaces-connected-sum}
 \Delta_\# P=\frac12\Delta_\#\mathcal W_{\mathrm{ten}}
             +\frac{\Delta_\# S}{12}+32\pi^2.
\end{equation}
Indeed, $\chi(M\# Z)=\chi(M)+\chi(Z)-2$. A strict Weyl
subadditivity theorem therefore is not, by itself, subadditivity for $P$.

Malchiodi--Malizia \cite[Theorem 1.1]{MalchiodiMalizia2025}
construct strict Weyl-energy competitors on the prescribed oriented
connected sum of two closed non-locally-conformally-flat four-manifolds,
provided they are not both self-dual or both anti-self-dual.
This applies to a complex-oriented complex-hyperbolic surface and
$\overline{\CP}^{\,2}$ with the reversed Fubini--Study metric:
the former is self-dual, the latter anti-self-dual, and neither is
locally conformally flat. Thus the theorem already supplies Weyl
competitors on the relevant blowup topology. Its small-neck gain does
not supply the joint scalar cost in
\eqref{eq:interfaces-fixed-topology}, or a uniform per-surgery saving
at the first-order scale of the balanced problem.

The revised locally-conformally-flat gluing theorem
\cite[Theorem 1.1]{MalchiodiMalizia2026} instead assumes a Bach-flat
factor that is neither self-dual nor anti-self-dual, and a positive-Yamabe
locally conformally flat factor not globally conformal to the round
sphere. A complex-hyperbolic factor is half-self-dual and does not
satisfy that theorem's hypothesis. The two results provide different
interaction mechanisms; neither may be substituted for the other's
assumptions. Both papers use $\mathcal W_{\mathrm{ten}}$.

For the original surgery family
\[
 M_j=X_j\# k_j\overline{\CP}^{\,2}\#\ell_j(S^1\times S^3),
 \quad A_j=c_1^2(X_j),\quad Q_j=k_j+4\ell_j,
 \quad h_j=Q_j/A_j,
\]
one can make the comparison independent of any chosen gluing scheme.
The classical recovery limits are
\[
 S_j^{\mathrm{ref}}=32\pi^2A_j,\qquad
 \mathcal W_{\mathrm{op},j}^{\mathrm{ref}}
       =\frac{4\pi^2}{3}A_j+12\pi^2k_j,
 \qquad \chi(M_j)=A_j/3+k_j-2\ell_j.
\]
Substitution in \eqref{eq:interfaces-joint-cost} gives, for every
smooth metric on $M_j$,
\begin{equation}\label{eq:interfaces-recovery-comparison}
 P(g)-8\pi^2Q_j
  =\frac{S(g)-S_j^{\mathrm{ref}}}{12}
   +2\bigl(\mathcal W_{\mathrm{op}}(g)
                 -\mathcal W_{\mathrm{op},j}^{\mathrm{ref}}\bigr).
\end{equation}
These reference values are simultaneous limits, not a claim that a
smooth metric realizes all of them. Formula
\eqref{eq:interfaces-recovery-comparison} isolates exactly what a
new interaction construction must improve. It does not subtract two
independently optimized functionals.

\subsection{Replication and first-order gain}
The effect of a cheaper cell depends on how many cells can actually be
replaced. Suppose $m_j$ blowup sites admit a fully matched construction
giving actual metrics with
\begin{equation}\label{eq:interfaces-replicated-cost}
 P(g_j)\le8\pi^2(Q_j-\delta m_j)+o(Q_j),
 \qquad 0<\delta<1,
 \qquad \frac{m_j}{k_j}\longrightarrow\theta\in(0,1].
\end{equation}
All transition and exterior errors are included in this hypothesis.
The balanced ratio is $k_j/Q_j\to2/9$. Theorem
\ref{thm:vg-first-order-reduction} therefore implies
\begin{equation}\label{eq:interfaces-replication-fraction}
 \limsup_j\frac{\inf_g\E(g)}{h_j}
     \le4-\frac89\theta\delta.
\end{equation}
To see this, divide the established relation
$\inf_g\E(g)=4\inf_gP(g)/(8\pi^2A_j)+O(h_j^2)$ by $h_j$,
insert \eqref{eq:interfaces-replicated-cost}, and use
$m_j/Q_j\to2\theta/9$. The bound $4-8\delta/9$ requires
$\theta=1$, or a saving $\delta$ already averaged over all blowup
sites. A saving tending to zero with the neck scale need not change
the leading coefficient.

This is a conditional consequence of an actual metric construction,
not the existence of such a construction. The all-metric monopole
supply on the original family supplies new solutions for a new metric;
it does not require a competitor to preserve the original glued
spinor. This leaves room for coupled curvature interactions while
keeping the same variational problem.

\subsection{Structure and realization retain their own roles}
The period constraints of Sections~\ref{sec:period-inertia} and
\ref{sec:joint-matrix}, the realization results of
Sections~\ref{sec:chern-boundary} and \ref{sec:virtual-realization},
and the slow-field and fixed-bundle constructions of
Sections~\ref{sec:slow-moduli} and \ref{sec:fixed-bundle-geometry}
describe geometric information beyond a numerical gap. Whole marked
fillings in Section~\ref{sec:whole-core} carry relative data not
determined by a terminal ALE space alone. Relating them to ideal
minimizing configurations \cite{Anderson2005} requires the relevant
topological matching and energy accounting. Such a relation may help
either a lower bound or a competing construction; it is not prescribed
as a prerequisite for every direct proof.

The density-independent endpoint and the balanced optimal coefficient
are precise tests of the present theory. They do not define the entire
space of actual suppliers or compatible geometries. A new supplier,
a stronger simultaneous-field constraint, or an exact realization is
an extension when it yields a conclusion beyond the already proved
range, with the required differential and global data supplied.
